\newcommand{\tipoRequerimento}{custo de disponibilidade}

\section{DOS FATOS E DOS DIREITOS}

\begin{enumerate}[resume]
\item A Unidade Consumidora \unidadeConsumidora{} enquadrada na modalidade de Iluminação Pública, vêm sendo faturadas com a cobrança de custo de disponibilidade, prática que afronta a regulamentação vigente e carece de amparo legal.
\item Nos termos do art. 290 da Resolução Normativa nº 1.000/2021, a distribuidora deve proceder ao faturamento da unidade consumidora enquadrada no Grupo B pelo maior valor entre o consumo de energia elétrica ativa e o custo de disponibilidade. Todavia, o § 1º do referido dispositivo estabelece, de forma expressa:
\Citacao{
    § 1º Não se aplica o custo de disponibilidade no faturamento de unidades consumidoras:

    I – da classe iluminação pública;

    II – atendidas por meio de sistemas isolados do tipo SIGFI ou MIGDI; e

    III – enquadradas na modalidade de pré-pagamento. (grifou-se)
}
\item No mesmo sentido, o art. 24, § 2º, da Resolução Normativa ANEEL nº 414/2010 é categórico ao estabelecer que:
\Citacao{
    Art. 24 O consumo mensal da energia elétrica destinada à iluminação pública deve ser apurado considerando as seguintes disposições:

    (...)

    § 2º Não se aplica a cobrança pelo custo de disponibilidade definida no art. 98 no faturamento individual de um ponto de iluminação pública. (grifou-se)
}\
\item Diante do exposto, conclui-se que o custo de disponibilidade não é aplicável às unidades consumidoras destinadas à iluminação pública, devendo o respectivo faturamento ocorrer exclusivamente com base no consumo efetivo de energia elétrica ativa.
\item Dessa forma, é incontroverso que a cobrança efetuada carece de respaldo na regulação setorial aplicável, impondo-se, por conseguinte, a sua imediata correção.
\item Para além da evidente afronta aos arts. 290, §1º, da REN 1.000/2021 e 24, §2º, da REN 414/2010, impõe-se ainda ressaltar o disposto no art. 323 da REN 1.000/2021, o qual assegura ao consumidor o direito à revisão do faturamento sempre que verificada a ocorrência de erro:
\Citacao{
    Art. 323. Constatado erro no faturamento, a distribuidora deve efetuar a revisão do faturamento, cobrando ou devolvendo ao consumidor a diferença apurada.

    (...)

    II - faturamento a maior: devolver ao consumidor e demais usuários, até o segundo ciclo de faturamento posterior à constatação, as quantias recebidas indevidamente nos últimos 60 ciclos de faturamento imediatamente anteriores à constatação. (Suspensos os efeitos referentes ao prazo de 60 (sessenta) ciclos, conforme DSP ANEEL 2.006, de 08.07.2024)
}
\item Assim, uma vez constatada a cobrança indevida, incumbe à distribuidora o dever de proceder à revisão e ao refaturamento, com a consequente restituição ao consumidor dos valores pagos a maior.
\item A manutenção de cobrança sabidamente ilegal pela concessionária afronta diretamente o princípio da legalidade (art. 37, caput, CF/88), além de caracterizar enriquecimento ilícito por parte da concessionária, em manifesta violação ao dever de observância da modicidade tarifária.
\end{enumerate}

\section{DOS PEDIDOS}

\begin{enumerate}[resume]
\item Ante o exposto, requer-se:
    \begin{enumerate}
    \item A imediata revisão do faturamento da Unidade Consumidora identificada, com a consequente exclusão da cobrança de custo de disponibilidade, nos termos do art. 290 da Resolução Normativa nº 1.000/2021, abrangendo o período correspondente aos últimos 10 (dez) anos, em conformidade com o disposto no Despacho nº 2.006/2024;
    \item Adequação definitiva do faturamento, de modo a observar estritamente o disposto no art. 290, § 1º, da Resolução Normativa 1.000/2021, bem como no art. 24, § 2º, da Resolução Normativa 414/2010, afastando a cobrança do custo de disponibilidade;
    \item A devolução em dobro dos valores indevidamente pagos, nos termos do art. 323, inciso II, da Resolução Normativa 1.000/2021;
    \item Proceda com a devolução de todos os valores cobrados a maior nos últimos 10 (dez) anos, conforme o Despacho nº 2.006 da Aneel, de 08 de julho de 2024, observando o prazo prescricional de 10 anos, conforme estabelecido no art. 205 do Código Civil. Conforme fundamentado na sentença proferida em 29 de setembro de 2023, pela 19ª Vara Cível Federal de São Paulo, no âmbito da Ação Civil Pública Cível nº 5024153- 93.2018.4.03.6100¹, no Ofício nº 00710/2024/PFANEEL/PGF/AGU², de 08 de maio de 2024 e no Ofício nº 00937/2024/PFANEEL/PGF/AGU, de 12 de junho de 2024³.
    \item Devolva os valores cobrados a maior em dobro, atualizados monetariamente pelo IPCA e juros de mora de 1\% ao mês calculados pro rata die, conformidade com o estabelecido na Resolução Normativa n° 1.000/2021 da Aneel.
    \item A expedição de resposta formal no prazo regulamentar, com a comprovação da adoção das providências cabíveis.
    \item Proceda com a devolução do valor integral, por meio de depósito bancário, na conta corrente do Município, conforme estabelece o § 7º, do art. 323 da REN n° 1.000/2021 da Aneel, comunicando por escrito o resultado do deferimento desta reclamação, por meio do endereço eletrônico: contato@halaconsultoria.com.
    \end{enumerate}
\end{enumerate}